%! Trigonometry and Polar Coordinates — Unit 7, Lesson 7.6 — Lesson Plan
\documentclass[10pt]{article}
\usepackage{precalculus-article}
\usepackage{precalculus-boxes}
\usepackage{graphicx}
\graphicspath{{images/}}
\newcommand{\TallMath}[1]{$\displaystyle #1\rule[-1.4em]{0pt}{3.2em}$}
\newcommand{\CourseName}{Precalculus}
\newcommand{\MeetingLength}{55 minutes}
\newcommand{\UnitNumberName}{Unit 7: Analytic Trigonometry and Polar Coordinates \quad}
\newcommand{\LessonNumberName}{Lesson 7.6: Trigonometry and Polar Coordinates}

\begin{document}

%----------------------------------------------------------------------
% Title Block
%----------------------------------------------------------------------
\begin{center}
  {\Large\bfseries \CourseName} \\[4pt]
  {\small\bfseries \UnitNumberName \LessonNumberName \quad (2--3 instructional periods, \MeetingLength\ each)}
\end{center}

%----------------------------------------------------------------------
% Primary Objective
%----------------------------------------------------------------------
\begin{tcolorbox}[colback=sky, colframe=navy, boxrule=0.9pt, arc=2mm,
    left=2mm, right=2mm, top=1.5mm, bottom=1.5mm]
  \textbf{Primary Objective:} Students will locate points in the plane using
  polar coordinates $(r,\theta)$, convert between polar and rectangular coordinate
  systems using $x = r\cos\theta$, $y = r\sin\theta$, $r = \sqrt{x^2+y^2}$, and
  $\theta = \arctan(y/x)$ (with quadrant adjustment), and represent complex numbers
  in both rectangular and polar form.
  \textit{(Big Idea: Functions)}
\end{tcolorbox}

\vspace{0.08in}

%----------------------------------------------------------------------
% Priority Ideas & Skills
%----------------------------------------------------------------------
\begin{skillbox}[Priority Ideas \& Skills]{goldbox}
\begin{tabularx}{\linewidth}{@{} p{0.46\linewidth} X @{}}
\textbf{AP Skills} & \textbf{Key Understandings (EKs)} \\
\midrule
\textbf{1.B} \textit{Express in equivalent forms} ---
  Convert a point between rectangular coordinates $(x,y)$ and polar
  coordinates $(r,\theta)$; express a complex number in both rectangular
  and polar form. Recognize that the same point has infinitely many polar
  representations.
&
\textbf{EK 3.13.A.1}: The polar coordinate system uses concentric circles
(radius $r$) and rays through the origin (angle $\theta$). The ordered pair
$(r,\theta)$ has $|r|$ as the circle radius and $\theta$ as the angle in
standard position. The same point has infinitely many polar representations,
e.g., $(r,\theta)$, $(r,\theta+2\pi)$, $(-r,\theta+\pi)$. \\[4pt]

\textbf{2.A} \textit{Identify information from representations} ---
  Given a polar coordinate pair, identify the circle and ray locating the
  point; given a rectangular point, determine polar coordinates; given a
  complex number, read its polar representation from the complex plane.
&
\textbf{EK 3.13.A.2}: Polar $\to$ rectangular: $x = r\cos\theta$,
$y = r\sin\theta$. \\[4pt]

&
\textbf{EK 3.13.A.3}: Rectangular $\to$ polar: $r = \sqrt{x^2+y^2}$;
$\theta = \arctan(y/x)$ for $x>0$; $\theta = \arctan(y/x)+\pi$ for $x<0$. \\[4pt]

&
\textbf{EK 3.13.A.4}: A complex number $a+bi$ corresponds to the point
$(a,b)$ in the complex plane. In polar form: $(r\cos\theta)+i(r\sin\theta)$
where $r = \sqrt{a^2+b^2}$ and $\theta = \arctan(b/a)$ (with quadrant
adjustment). \\
\end{tabularx}
\end{skillbox}

\vspace{0.08in}

%----------------------------------------------------------------------
% Vocabulary
%----------------------------------------------------------------------
\begin{skillbox}[{Vocabulary, Concepts, \& Theorems}]{greenbox}
\begin{tabularx}{\linewidth}{@{} p{0.26\linewidth} X @{}}
\textbf{Term} & \textbf{Definition / Note} \\
\midrule
polar coordinate system &
  A coordinate system based on concentric circles centered at the origin
  and rays from the origin. Points are located by $(r,\theta)$. \\[4pt]
pole &
  The fixed reference point (the origin) in the polar coordinate system. \\[4pt]
polar axis &
  The reference ray from the pole in the direction of positive $x$. \\[4pt]
polar coordinates $(r,\theta)$ &
  Ordered pair where $|r|$ is the distance from the pole and $\theta$ is
  the angle (in standard position) whose terminal ray passes through the point. \\[4pt]
multiple representations &
  $(r,\theta)$ also equals $(r,\theta+2\pi k)$ for integer $k$, or $(-r,\theta+\pi)$. \\[4pt]
complex plane &
  Coordinate plane where the horizontal axis gives the real part and the
  vertical axis gives the imaginary part of a complex number. \\[4pt]
polar form of a complex number &
  $(r\cos\theta)+i(r\sin\theta)$, where $r = \sqrt{a^2+b^2}$ and $\theta$ is the
  argument of $a+bi$. \\
\end{tabularx}
\end{skillbox}

\vspace{0.08in}

%----------------------------------------------------------------------
% Activate Prior Knowledge
%----------------------------------------------------------------------
\begin{skillbox}[Activate Prior Knowledge \& Spiral Review]{sky}
\begin{tabularx}{\linewidth}{@{}X c@{}}
\begin{minipage}[t]{0.96\linewidth}\vspace{0pt}
\textbf{Spiral topics activated during Warm-Up (first 8--10 minutes):}
\begin{itemize}[leftmargin=1.4em, itemsep=2pt]
  \item \textbf{Distance from the origin}: $r = \sqrt{x^2+y^2}$ generalizes
    the distance formula from Unit 1 --- fluency needed to convert rectangular
    to polar.
  \item \textbf{Unit-circle cosine and sine}: Evaluating $\cos\theta$ and
    $\sin\theta$ at standard angles directly drives the polar$\to$rectangular
    formulas.
  \item \textbf{Inverse tangent}: $\arctan(y/x)$ gives the reference angle;
    students must recall from Lessons 3.9--3.11 that a quadrant check is
    required to find $\theta$.
  \item \textbf{Complex numbers (Algebra 2 prerequisite)}: The rectangular form
    $a+bi$ and the complex plane provide context for EK 3.13.A.4.
\end{itemize}
\end{minipage}
&
\end{tabularx}
\end{skillbox}

\vspace{0.08in}

%----------------------------------------------------------------------
% Hook
%----------------------------------------------------------------------
\begin{skillbox}[Hook --- Entry Event]{sky}
\textbf{Scenario:} Write the point $(1,1)$ in rectangular form on the board.

\textbf{Discussion prompts:}
\begin{enumerate}[leftmargin=1.4em, itemsep=3pt]
  \item \textit{``You know $(1,1)$ means 1 unit right and 1 unit up from the
    origin. What is the \emph{distance} of this point from the origin?''}
    (Students: $\sqrt{1^2+1^2} = \sqrt{2}$.)
  \item \textit{``What \emph{angle} does the segment from the origin to $(1,1)$
    make with the positive $x$-axis?''} (Students: $\pi/4$, or $45°$.)
  \item \textit{``Could you describe the location of $(1,1)$ using just a
    distance and an angle instead of $x$ and $y$?''}
\end{enumerate}

\textbf{Key tension to surface:} The rectangular system measures along two
perpendicular axes. The polar system describes position by distance from a
central point and an angle of rotation. Both locate the same point --- neither
is more ``correct.'' Radar, sonar, and navigation instruments naturally output
a distance and a bearing, i.e., polar coordinates.

\textbf{Transition:} ``Today we'll define the polar coordinate system formally,
develop the conversion formulas between polar and rectangular, and see how
complex numbers connect the two systems.''
\end{skillbox}

\vspace{0.08in}

%----------------------------------------------------------------------
% Lesson
%----------------------------------------------------------------------
\begin{skillbox}[Lesson --- Instruction Sequence]{sky}
\begin{multicols}{2}

\textbf{Step 1 --- The Polar Coordinate System (EK 3.13.A.1)} (10 min)

Define the pole (origin) and polar axis (positive $x$-direction). Explain
$(r,\theta)$: $|r|$ is the radius of the circle on which the point lies;
$\theta$ is the angle in standard position whose terminal ray contains the point.
Together plot $(2,\pi/6)$, $(3,-\pi/4)$, and $(-2,\pi/3)$. Emphasize: negative
$r$ reflects through the pole. Introduce multiple representations:
$(2,\pi/6) = (2,\pi/6+2\pi) = (-2,\pi/6+\pi)$.

\columnbreak

\textbf{Step 2 --- Polar $\to$ Rectangular (EK 3.13.A.2)} (10 min)

Derive $x = r\cos\theta$, $y = r\sin\theta$ from scaling the unit-circle point
$(\cos\theta, \sin\theta)$ by $r$.
Worked example: $(4,\pi/3)$ gives
$x = 4\cos(\pi/3) = 4 \cdot \tfrac{1}{2} = 2$ and
$y = 4\sin(\pi/3) = 4 \cdot \tfrac{\sqrt{3}}{2} = 2\sqrt{3}$.
Students practice converting $(3, 3\pi/4)$ in guided notes.

\end{multicols}

\begin{multicols}{2}

\textbf{Step 3 --- Rectangular $\to$ Polar (EK 3.13.A.3)} (12 min)

Display $r = \sqrt{x^2+y^2}$ and the $\arctan$ formula. Stress the quadrant
adjustment: $\arctan(y/x) \in (-\pi/2,\pi/2)$; add $\pi$ if $x < 0$.
Q1 example: $(3,3)$ gives $r = 3\sqrt{2}$, $\theta = \pi/4$.
Q2 example: $(-\sqrt{3},1)$ gives $r = 2$,
$\theta = \arctan(1/(-\sqrt{3}))+\pi = -\pi/6+\pi = 5\pi/6$.
Students practice $(-1,-1)$ in guided notes (Q3: add $\pi$).

\columnbreak

\textbf{Step 4 --- Complex Numbers in Polar Form (EK 3.13.A.4)} (10 min)

Map $a+bi$ to $(a,b)$ in the complex plane. Identify $r = \sqrt{a^2+b^2}$
(modulus) and $\theta = \arctan(b/a)$ with quadrant check (argument).
Polar form: $(r\cos\theta)+i(r\sin\theta)$.
Worked example: $1+i\sqrt{3}$: $r = 2$, $\theta = \pi/3$;
polar form $= 2\cos(\pi/3)+2i\sin(\pi/3)$.

\end{multicols}

\smallskip
\textbf{Pacing note:} Steps 1--2 in Period 1; Steps 3--4 and group activity
in Period 2; review and exit ticket in Period 3 if needed.
\end{skillbox}

\vspace{0.08in}

%----------------------------------------------------------------------
% Active Monitoring
%----------------------------------------------------------------------
\begin{skillbox}[Active Monitoring]{redbox}
\begin{itemize}[leftmargin=1.4em, itemsep=3pt]
  \item \textbf{Quadrant error in $\arctan$}: Students apply $\arctan(y/x)$
    without checking the quadrant. Prompt: \textit{``Plot the point first.
    Which quadrant is it in? Does your $\arctan$ result land there?''}
  \item \textbf{Sign of $r$}: Students think $r$ must be positive. Address:
    $(-r,\theta)$ represents the same point as $(r,\theta+\pi)$; both are valid.
    For rectangular$\to$polar we typically take $r > 0$ and adjust $\theta$.
  \item \textbf{Cosine/sine swap}: Students compute $x = r\sin\theta$ or
    $y = r\cos\theta$. Prompt: ``Cosine gives the \emph{horizontal} component;
    sine gives the \emph{vertical} component.''
  \item \textbf{Uniqueness confusion}: Students expect one polar form. Clarify:
    rectangular coordinates are unique; polar coordinates are not. Ask students
    to list at least two valid polar representations of the same point.
\end{itemize}
\end{skillbox}

\vspace{0.08in}

%----------------------------------------------------------------------
% Group Work & Differentiation
%----------------------------------------------------------------------
\begin{skillbox}[Group Work \& Differentiation]{redbox}
Students work in groups of 3--4 on the tiered activity (teacher-assigned or
student self-selected with guidance).

\begin{itemize}[leftmargin=1.4em, itemsep=4pt]
  \item \textbf{Tier R (Reinforce)}: Given four polar coordinate pairs
    (unit-circle angles), complete a step-by-step table converting to
    rectangular coordinates using $x = r\cos\theta$ and $y = r\sin\theta$.
  \item \textbf{Tier A (Apply)}: Convert three rectangular points to polar
    form; find two additional polar representations of a given point;
    convert one complex number to polar form.
  \item \textbf{Tier E (Extend)}: Write a general formula for all polar
    representations of a Q2 point; show algebraically that $(-r,\theta)$ and
    $(r,\theta+\pi)$ name the same point; reason about why polar representations
    are not unique while rectangular ones are.
\end{itemize}
\end{skillbox}

\vspace{0.08in}

%----------------------------------------------------------------------
% Individual Work & Assessment
%----------------------------------------------------------------------
\begin{skillbox}[Individual Work \& Assessment]{redbox}
Exit ticket administered independently (final 5--7 minutes of Period 2 or 3).

\begin{enumerate}[leftmargin=1.4em, itemsep=4pt]
  \item Convert the polar coordinates $(4,\pi/3)$ to rectangular coordinates.
  \item Convert the rectangular coordinates $(-3,3)$ to polar form ($r > 0$,
    $0 \le \theta < 2\pi$).
  \item Write the complex number $2-2i$ in polar form.
\end{enumerate}

\textbf{Data use:} Sort exit tickets by error type --- quadrant errors in
$\arctan$, sign-of-$r$ errors, cosine/sine swap errors. Use distribution
to determine whether a brief re-teach opener is needed at Lesson 3.14.
\end{skillbox}

\vspace{0.08in}

%----------------------------------------------------------------------
% Reinforcement & Extension
%----------------------------------------------------------------------
\begin{skillbox}[Reinforcement \& Extension]{goldbox}
\textbf{Homework (Lesson 7.6):} 5--6 problems covering polar$\to$rectangular
conversions, rectangular$\to$polar conversions, multiple representations of a
point, one complex number problem, and 2 AP-style MC items.

\textbf{Preview of Lesson 3.14 --- Polar Function Graphs:}
Ask students: \textit{``Now that you can locate any point using $(r,\theta)$,
what happens when $r$ is a \emph{function} of $\theta$? In Lesson 3.14 we'll
graph polar equations like $r = 1+\cos\theta$ and discover curves that are
difficult to describe in rectangular form.''}
\end{skillbox}

\end{document}
